% ============================================================================
% PART II, CH C — INTERPRETABILITY TUNING       [slug: interpretability-tuning]
% Third of the three Part II chapters (split 2026-07-14): the METHOD — "a
% whole framework around interpretability knobs and metrics" (Matteo), plus
% ALL of Part II's experiments (pooled here, none in B).
% RESTRUCTURED 2026-07-14 (Matteo, second pass — the 10<->11 connection):
% the METHOD comes first and the TWO ROUTES are its final DISPATCH step, not
% a preceding philosophy; the previous chapter's taxonomy is the LOOKUP
% TABLE the method consults (metrics + adjustment classes). A demonstration
% section is mandatory.
% GOAL FRAMING (Matteo): Part I already delivered interpretable models —
% this chapter works toward TUNABLY interpretable / MORE interpretable
% models. (Echoes SQ-3 title-ladder D "Tunably Interpretable
% Recommendation".)
% RASHOMON here = implicit base assumption for Route 1, one sentence, never
% measured (scoping in chapter A's header).
% STATUS: methodology derived (2026-07-12_14xx quintet + 2026-07-13_2326);
% experiments NOT designed yet.
% Material map: ../outline.md#interpretability-tuning
% ============================================================================
\cleardoublepage
\chapter{Interpretability Tuning (7--9)}   % page goal — scaffolding, DELETE before submission
\label{ch:iq-tuning}

% CHAPTER INTRO (bonbon, unnumbered): the goal shift — Part I delivered
% interpretable models; this chapter is about making them TUNABLY / MORE
% interpretable. The method in one sentence: decompose the explanation's
% mathematics into every interpretability-affecting component, then steer
% each component through the channel its nature admits.

\section{The Method: Explanation-Backward Tracing}
\label{sec:iq-method}
\vault{2026-07-13_2326_part-ii-framework-knob-entry-points-explanation-backward-tracing, 2026-07-14_1218_descriptor-gap-always-a-gap-fidelity-vs-legibility}
% The executable pipeline (Matteo's framework, vault 2026-07-13_2326 +
% second-pass refinement 2026-07-14 — status "decent starting point, far
% from worked out"; deliberately tailored to our model, generality = Future
% Work).
\begin{itemize}
    \item Fix the explanation object — \emph{the explanation is a formula}. \vault{2026-07-13_2326_part-ii-framework-knob-entry-points-explanation-backward-tracing}
    \item Quantify interpretability on it via the taxonomy's metrics — sparsity first.
    \item Backtrace: decompose the mathematics into \emph{every single component} affecting those metrics.
    \item Per component: desired direction + adjustment class (looked up in the taxonomy).
    \item Lastly: dispatch each component to its channel (next section).
\end{itemize}

\section{Tunably More Interpretable Models through Two Channels}
\label{sec:iq-two-routes}
\label{sec:iq-uw-knob}   % UW feature use folded in as a bullet (2026-07-14)
\vault{2026-07-12_1409_two-route-rashomon-methodology, 2026-07-12_1410_knob-metric-axis-and-principledness-spectrum}
% Retitled 2026-07-14 (Matteo; was "Two Routes, Working Together") —
% "channel" replaces "route" in prose-facing text; "Route 1/Route 2" stays
% below as the established vault shorthand (2026-07-12_14xx entries) —
% whether the shorthand also becomes "Channel 1/2" is a prose-time
% micro-decision.
% The dispatch — the method's final step, not a rival-philosophies debate
% (vault 2026-07-12_1409 + Matteo's second pass 2026-07-14):
\begin{itemize}
    \item \emph{Model-structural} knobs $\to$ Route 1: search within the accuracy window $\varepsilon$ (Rashomon set = implicit base assumption, never measured). \vault{2026-07-12_1409_two-route-rashomon-methodology}
    \item \emph{Differentiable} knobs $\to$ Route 2: adopted into the training loss as surrogate terms. \vault{2026-07-12_1410_knob-metric-axis-and-principledness-spectrum}
    \item \emph{Validation-type} knobs $\to$ regularizers and/or validation numbers feeding Route-1 selection.
    \item The channels work together: scalarizing (Route 2) is picking a point on the front Route 1 explores. \vault{2026-07-12_1409_two-route-rashomon-methodology}
    \item Both channels in the thesis; Route 2 emphasized; bare-regularizer landing acceptable. \vault{2026-07-12_1409_two-route-rashomon-methodology}
    \item Understandability-weighted feature use — our own mechanism, dual-natured: $(1{-}s_f)$-weighted group-lasso in the loss, or up-front feature selection. \vault{2026-07-12_1413_understandability-weighted-feature-use-and-sparsity-axes, 2026-06-11_0905_r7-user-perceivable-feature-vocabulary}
    \item UW sparsity is biased toward user-legible features — operationalizes R7. \vault{2026-06-11_0905_r7-user-perceivable-feature-vocabulary}
    \item Novelty boundary carried exactly: Rudin's principle, our mechanism; $s_f$ source parked. \vault{2026-07-12_1413_understandability-weighted-feature-use-and-sparsity-axes}
\end{itemize}
% BLOCK — the principledness spectrum (vault 2026-07-12_1410, idea 2), the
% quality ladder for whatever enters the loss: ad-hoc penalty -> penalty
% with probabilistic reading (MAP) -> fully derived. "Just another
% regularizer" = a documented level-1 position, not a failure; future work =
% lifting properties up the spectrum. Facts: the host's entropy-based
% regularizers are already ~level 2; the one clear level-1 outlier is the
% separation hinge -> the DPP lift is the obvious first demonstrated climb.

\section{Demonstrating the Method}
\label{sec:iq-demo}
\label{sec:iq-experiments}   % experiments folded in here (Matteo 2026-07-14
\vault{2026-07-13_2326_part-ii-framework-knob-entry-points-explanation-backward-tracing, 2026-06-16_2014_popularity-bias-prototype-interpretability-ablation}
                             % — the demonstration IS the experiment program,
                             % no separate title)
% MANDATORY section (Matteo, 2026-07-14). Two movements, one arc:
% MOVEMENT 1 — the analytic walk: fix the fI/fU score formula -> sparsity as
% the worked interpretability function -> backtrace to the concrete
% components (profiles, prototype matrix, feature blocks, regularizer
% terms) -> desired directions -> adjustment classes from the taxonomy ->
% dispatch table as the payoff artifact (which knob goes to which channel,
% and why). The four sparsity axes and the understandability-weighted
% mechanism (bulleted in the channels section) fall out of this walk as
% outputs, demonstrating the method produces knobs rather than assuming them.
% MOVEMENT 2 — the empirical proof [NOT YET DESIGNED; ALL Part II
% experiments live here]:
%  1. Measuring the lineage: facet instrumentation on the trained Part I
%     checkpoints — profile entropy/sharpness, feature-explained fraction,
%     purity, determinacy (the taxonomy's metrics computed on our own models).
%  2. Route-1 studies: knob search under the accuracy-tolerance window
%     (the regularizer weights as knobs, selection on measured metrics).
%  3. Route-2 studies: the understandability-weighted group-lasso; the DPP
%     separation lift as the principledness climb.
%  4. Accuracy-cost numbers (the part-opening chapter's inoculation passage
%     quotes these): host comparisons fI/fU/fUfI vs their hosts.
%  5. The honest unknowns AS experiment design (vault 2026-07-13_2326):
%     weight selection (epsilon protocol), gradient conflict between
%     objectives, surrogate gap (proxy moves, true metric doesn't). All
%     three tractable as experiments where the theory is hard.
% Reality check carried from the founding notes: the host already trains a
% compound loss (rec + two weighted regularizers) — adding a term is
% mechanically routine; the three unknowns above are the real work.
